\section{Tame reducible residual inertia at the first wild index}\label{sec:tame-reducible}
This section recovers the character at $p=5$ and bounds the possible wild image of an unequal-character comparison at the general first wild index. The first reduction identifies the canonical lines from a small common point field; two opposite quotient actions, diagonal visibility, and a cyclotomic quotient obstruction force seven matrix dimensions when $p\equiv1\pmod4$.

\begin{proposition}\label{prop:square-determinant}
At $e=p(p-1)$, $p\ge5$, a $p^3$-torsion comparison between two potentially supersingular height-two groups with a common square finite determinant determines their special characters.
\end{proposition}
\begin{proof}
Twist both groups by the inverse of a square root of their finite determinant. The finite special inertia is now norm one in $\F_{p^2}^\times$, so its good degree $m$ divides $p+1$. The good field has absolute index $a=em\le p(p^2-1)$.

The proof of Lemma~\ref{lem:weighted-kernel} without the tame weight congruence gives the recurrence
\[
                      q_0=1,\qquad q_{j+1}=\min\{a+q_j,1+pq_j\}.
\]
Indeed a degree-$p$ coefficient is still divisible by the uniformizer, and the rest of the logarithm and integral Cartier argument is unchanged. Here $q_1=p+1$ and
\[
                q_2=\min\{a+p+1,p^2+p+1\}\ge a/(p-1).
\]
Apply the graph-closure proof of Theorem~\ref{thm:common-determinant}. If the special characters differed, its target would have product-height-two special fibre, and its special kernel would be $\ker F^3$. The unweighted bound would kill that kernel by $F^2$, which is impossible.
\end{proof}

We recall the two-slope notation over the degree-$N$ good field. Put $s=(p+1)e$. The high line $H$ and the low points have valuations
\begin{equation}\label{eq:two-slope-B}
        \alpha=s-B,\qquad\beta=B/p,\qquad0<B<pe.
\end{equation}
The high line is rational over $L$, so $B$ is an integer; its tame character is $\zeta^{c-B}$, where $c$ is the special tangent exponent. These formulas follow from the vertices $(1,eN),(p,(p-1)B),(p^2,0)$ of the Newton polygon. The high-line character follows by multiplying the cotangent factor $\zeta^c$ by the leading coefficient factor $\zeta^\alpha$, noting that $s$ is divisible by $N$ when $e=p(p-1)$.

\begin{lemma}\label{lem:one-prime-parameter}
Suppose that two two-slope groups have rational $p$-torsion over $L$ and isomorphic $p^2$-torsion. Write $B_i=pb_i$, $0<b_i<e$. If $p\nmid b_1$, the high lines agree, $b_1=b_2$, and the special characters agree.
\end{lemma}
\begin{proof}
Take a residual vector outside both high lines and corresponding second-level lifts. Their valuations are $b_i/q$. Their common point field is Galois of degree at most $q$, since its conjugates are rational first-level translates. The denominator of $b_1/q$ forces degree $q$, so every translation occurs.

In this field the coordinate values are $b_i$. For every nonidentity $g$, the first model gives the exact break equality $v(gx_1-x_1)=b_1+b(g)$; the second gives the inequality $v(gx_2-x_2)\ge b_2+b(g)$. If the high lines differed, a translation high in the first and low in the second would imply
\[
                      qb_2\ge b_2+q(s-pb_1)-b_1.
\]
The difference of the left side and right side is increasing in both $b_i$, and at $b_1=b_2=e-1$ it is $-q(p+1)<0$. This is impossible. With the high lines identified, a low translation gives $(q-1)(b_2-b_1)\ge0$, and a high translation gives $(p^3+1)(b_1-b_2)\ge0$. Thus $b_1=b_2$. The equal high-line characters then give equal $c_i$ modulo $N$.
\end{proof}

\begin{corollary}\label{cor:first-wild-tame-pattern}
Let $p\ge5$, $e_K=p(p-1)$ and $\mu_p\subset K$. Suppose two potentially supersingular height-two groups have isomorphic $p^3$-torsion, tame reducible residual inertia and different special characters. Over the common tame good field of degree $p^2-1$, both first-level polygons have two slopes and parameters
\[
                  B_i=p^2k_i,\qquad1\le k_i\le p-2.
\]
Their high lines are the two distinct residual tame eigenspaces and are exchanged under the comparison.
\end{corollary}
\begin{proof}
Proposition~\ref{prop:square-determinant} excludes a square finite determinant. A reducible one-slope leading-coordinate action is scalar, so its determinant is square. Thus both models have two slopes. Their residual modules are rational over $L$, and Lemma~\ref{lem:one-prime-parameter} implies that both parameters $b_i=B_i/p$ are divisible by $p$. The strict range $0<B_i<pe$ gives the displayed bounds on $k_i$.

If the high lines coincided, equality of their characters would give $c_1-c_2\equiv B_1-B_2\pmod{p^2-1}$. Equality of the finite norm characters gives $c_1\equiv c_2\pmod{p-1}$. Hence $k_1\equiv k_2\pmod{p-1}$, and their range forces equality; the high-line character then gives $c_1=c_2\pmod{p^2-1}$, a contradiction. The two lines are distinct. A scalar residual action would again have square determinant, so these are the two distinct tame eigenspaces.
\end{proof}

For the next three lemmas, $K$ is complete discretely valued of mixed characteristic $(0,p)$ with algebraically closed residue field, $p\ge5$, $\mu_p\subset K$, and $e_K=p(p-1)$. The field $L/K$ is the tame extension of degree $N=p^2-1$. For an actual comparison whose $p$-torsion is rational over $L$, write $W\subset1+pM_2(\Z/p^3)$ for its common third-level image. Its first space $W_1$ is the image of $1+pA\mapsto A\pmod p$, and its terminal space is $W_2=\{B:1+p^2B\in W\}$.

\begin{lemma}\label{lem:first-layer-diagonal-visibility}
Let $p\ge5$ and $e_K=p(p-1)$ in the stated base setting. Suppose two whole potentially supersingular height-two groups have isomorphic actual $p^3$-torsion, rational $p$-torsion over $L$, parameters
\[
 B_i=p^2k_i,\qquad1\le k_i\le p-2,
 \qquad H_1=\F_pe_1,\quad H_2=\F_pe_2.
\]
For the common actual first congruence space $W_1$, both functions $A\mapsto A_{11}$ and $A\mapsto A_{22}$ are nonzero. No bound on the full division-field degree is assumed.
\end{lemma}
\begin{proof}
Suppose $A_{11}=0$ for every $A\in W_1$, and choose a second-level point $Q$ with $pQ=e_1$. Its orbit consists of translations by $W_1e_1\subset\F_pe_2$, so its point field has degree at most $p$. In the second model its coordinate has valuation $k_2/p$, whose denominator forces degree exactly $p$. Its stabilizer is the kernel of the column map on the abelian first-layer image; hence the field is Galois and its nonidentity automorphisms translate by nonzero multiples of $e_2$.

Since $B_i=p^2k_i>e$, a division of the high point in the first model has valuation $(N-pk_1)/p$, while the second model gives $k_2/p$. On the common degree-$p$ field both normalized values, $N-pk_1$ and $k_2$, are prime to $p$. The translating point is low in the first model and high in the second. The exact prime-valuation break formula therefore assigns the same automorphism the two breaks
\[
 p(p+1)k_1-N,\qquad p^2N-(p^3+1)k_2.
\]
Equality would require
\[
 p(p+1)k_1+(p^3+1)k_2=p^4-1.
\]
The left side is at most $(p-2)(p^3+p^2+p+1)$, which is smaller than $p^4-1$ by $(p+1)(p^2+1)$. This is impossible. Interchanging the models and axes proves that the $22$-projection is also nonzero.
\end{proof}

\begin{lemma}\label{lem:first-wild-cyclotomic-break}
Let $p\ge5$, $e_K=p(p-1)$, $\mu_p\subset K$, and $[L:K]=N=p^2-1$ as above. If $\mu_{p^2}\subset L$, then $\mu_{p^2}\subset K$, the extension $C=L(\mu_{p^3})/L$ is cyclic of degree $p$, and every nonidentity element has lower break $N^2$.
\end{lemma}
\begin{proof}
The degree of $K(\mu_{p^2})/K$ divides both $p$ and the prime-to-$p$ integer $N$, so this extension is trivial.
Choose a Cohen coefficient ring $\Lambda\subset\mathcal O_K$ with the same perfect residue field, using \citet[Theorem~10.160.8, Tag~032A]{StacksCohen}, and put $K_0=\operatorname{Frac}\Lambda$. The shifted polynomial $\Phi_{p^2}(1+T)$ is Eisenstein of degree $p(p-1)=e_K$. Its root $\zeta_{p^2}-1$ is a uniformizer of $K$, and successive expansion in that uniformizer with coefficients in $\Lambda$ gives $K=K_0(\zeta_{p^2})$. The analogous $p^3$-cyclotomic polynomial gives a cyclic degree-$p$ extension $K_3/K$ with uniformizer $\pi_3=\zeta_{p^3}-1$. Every nonidentity automorphism satisfies
\[
 v_{K_3}(\sigma\pi_3-\pi_3)=v_{K_3}(\zeta_p-1)=p^2,
\]
so its lower break is $N$. The compositum $C=LK_3$ has degree $p$ over $L$. In its normalized valuation, $v_C(\pi_3)=N$ and $v_C(\pi_L)=p$. Choose integers $a,b$ with $aN+bp=1$. Then $\pi_3^a\pi_L^b$ is a uniformizer, $p\nmid a$, and its nonidentity displacement ratio has valuation $N^2$. Thus
\begin{equation}\label{eq:general-cyclotomic-last-break}
                   b(C/L)=N^2.
\end{equation}

\end{proof}

\begin{lemma}\label{lem:traceless-first-wild-exclusion}
Under the hypotheses of Lemma~\ref{lem:first-layer-diagonal-visibility}, assume the residual inertia is tame reducible, the finite determinant is nonsquare, and the ratio of the two residual eigencharacters differs from one. The two actual congruence layers cannot both be $\mathfrak{sl}_2(\F_p)$ in the high-axis basis. This assertion holds for every $p\ge5$.
\end{lemma}
\begin{proof}
Suppose $W_1=W_2=\mathfrak{sl}_2(\F_p)$. On the good field the determinant is cyclotomic. Its first layer is traceless, so $\mu_{p^2}\subset L$. Lemma~\ref{lem:first-wild-cyclotomic-break} gives a nontrivial degree-$p$ cyclotomic quotient $C/L$ with lower break $N^2$.
On the terminal layer, $\det(1+p^2B)=1+p^2\operatorname{tr}B=1$. The nonzero cyclotomic character modulo $p^3$ consequently descends to a nonzero linear functional $\ell$ on $W_1=\mathfrak{sl}_2(\F_p)$. It is invariant under tame conjugation, so it vanishes on the nontrivial off-diagonal weights and is a nonzero multiple of $A_{11}$.

Take a second-level point $Q$ with $pQ=e_1$ and put $D=L(Q)$. Evaluation of $\mathfrak{sl}_2$ on $e_1$ is onto $\F_p^2$, so $D/L$ is Galois of degree $p^2$. Its point stabilizer maps to $\F_pE_{12}$ in $W_1$; the terminal layer and this off-diagonal line are both killed by the determinant. Thus the entire point stabilizer is killed by $\ell$, and $C\subset D$. In the second model $Q$ has normalized $D$-valuation $pk_2$. Every lift of a nonidentity element of $\operatorname{Gal}(C/L)$ translates $Q$ by a point with nonzero $e_1$ coefficient, which is low in that model. Its displacement has valuation $p^3k_2$. The general ramification inequality gives
\[
          b_D(g)\le p^3k_2-pk_2=pNk_2\le pN(p-2)<N^2.
\]
For a nonidentity quotient element, the norm-polynomial formula averages the indices $b_D(g)+1$ over its fibre; see \citet[Proposition~3.7 and its full proof, pp.~14--15]{ColmezLocal}. It therefore gives $b_C<N^2$, contradicting~\eqref{eq:general-cyclotomic-last-break}. This formula assumes characteristic zero and perfect residue, and applies here without a finite-residue hypothesis.
The contradiction proves the lemma.
\end{proof}

\begin{proposition}\label{prop:opposite-quotient-image}
Under the hypotheses of Corollary~\ref{cor:first-wild-tame-pattern}, suppose $p\equiv1\pmod4$. The wild image over $L$ on the $p^3$-torsion has order at least $p^7$.
\end{proposition}
\begin{proof}

First quotient both groups by either of the two common stable residual lines. These quotients have isomorphic $p^2$-torsion: for a line $J$ the relevant group is
\[
 (M/J)[p^2]=\{x\in M[p^3]:p^2x\in J\}/J.
\]
One quotient is canonical and the other noncanonical. Since $B_i=p^2k_i>e$, their parameters are respectively
\[
       p(p^2-1-pk_i),\qquad pk_j.
\]
These are the quotient formulas of Proposition~\ref{rc-prop:quotients}, which do not require $p\nmid e$. After division by $p$, both integers are prime to $p$, and
\[
               p^2-1-pk_i\ge2p-1>p-2\ge k_j.
\]
If the common residual quotient representation were trivial over $L$, Lemma~\ref{lem:one-prime-parameter} would equate these unequal integers. Its image is a $p$-group, being a subquotient of the original third-level image over $L$, so nontriviality means nontrivial wild inertia. The same conclusion holds for the opposite quotient.

Write $P$ for the original wild image, and choose a basis on the two residual tame lines. The first-level module is rational over $L$, so $P\subset1+pM_2(\Z/p^3)$. Its first and terminal matrix spaces $W_1,W_2\subset M_2(\F_p)$ satisfy
\[
 W_1\subset W_2,\qquad[W_1,W_1]\subset W_2,
 \qquad |P|=p^{\dim W_1+\dim W_2}.
\]
Indeed $(1+pA+p^2B)^p\equiv1+p^2A\pmod{p^3}$, and a group commutator reduces to $1+p^2[A,B]$.

The quotient by the first residual axis replaces a Tate lattice $T$ by $T+\Z_p(e_1/p)$. In the basis $(e_1/p,e_2)$ a matrix $\rho$ becomes
\[
 \begin{pmatrix}\rho_{11}&p\rho_{12}\\
                 \rho_{21}/p&\rho_{22}\end{pmatrix}.
\]
For $\rho=1+pA+p^2B$, its residual lower coefficient is $A_{21}$. The nontrivial residual wild action of this quotient makes the $21$-projection of $W_1$ nonzero. The opposite quotient makes its $12$-projection nonzero.

Let $a,d\in\F_p^*$ be the residual tame eigenvalues. Their product is the nonsquare finite determinant, because $\chi_p\equiv1\pmod p$. The ratio $r=a/d=a^2/(ad)$ is nonsquare. Since $p\equiv1\pmod4$, both $1$ and $-1$ are squares, and the three adjoint weights $1,r,r^{-1}$ are distinct. Tame spectral projectors preserve $W_1$ and therefore put both $E_{12}$ and $E_{21}$ in it. They also belong to $W_2$, as does their commutator $E_{11}-E_{22}$. These three matrices are independent. Consequently $\dim W_1\ge2$ and $\dim W_2\ge3$. This is the first, five-dimensional bound.

Lemma~\ref{lem:first-layer-diagonal-visibility} makes both diagonal projections nonzero. Tame projectors isolate a nonzero diagonal direction, independent of $E_{12},E_{21}$. Thus $\dim W_1\ge3$, and the power inclusion gives $\dim W_2\ge3$. If their total dimension were six, they would coincide. The two off-diagonal units and their commutator would then give
\[
                    W_1=W_2=\mathfrak{sl}_2(\F_p),
\]
contrary to Lemma~\ref{lem:traceless-first-wild-exclusion}. The total dimension is at least seven, proving the proposition.
\end{proof}

The congruence condition separates the opposite tame weights. When $p\equiv3\pmod4$, a nonsquare ratio can be $-1$; two nonzero off-diagonal projections then need not supply the two individual matrix units. The proposition makes no assertion in that case and constructs no unequal-character pair.

\begin{lemma}\label{lem:opposite-line-conductor}
Let $p$ be a prime, and let $T/S$ be a totally ramified Galois extension of complete discretely valued mixed-characteristic fields with the same perfect residue field, group $\F_p^2$, and two lower breaks $t<u$, with last line $V$. Suppose a coordinate $x$ has $v_p(A)=1$, where $A=v_T(x)$, and displacement $D_1$ on a line $H\setminus\{0\}$ and $D_0$ off it, with $A<D_0<D_1<v_T(p)$. If $H\ne V$, the last upper break is $(D_0-A)/p$.
\end{lemma}
\begin{proof}
Use the expansion~\eqref{eq:first-change}. A nonzero element of $H$ has break $t$ and displacement $D_1$, while a nonzero element of $V$ has break $u$ and displacement $D_0$. Both possible first-change degrees increase strictly with the break. Thus the first change at $t$ would occur before $D_0$ unless it cancels. Cancellation forces a nonzero ordinary term of degree $J=A+(p-1)t$. At $u$ that term is uniquely first, giving $J+u=D_0$. The last upper break is therefore $(u+(p-1)t)/p=(D_0-A)/p$.
\end{proof}

\begin{theorem}\label{thm:tame-five-recovery}
Let $K$ have algebraically closed residue field of characteristic five, $e_K=20$, and $\mu_5\subset K$. Two potentially supersingular height-two groups with isomorphic $125$-torsion and tame reducible residual inertia have equal finite special characters.
\end{theorem}
\begin{proof}
Assume otherwise. The determinant comparison fixes the common finite determinant; a square is excluded by Proposition~\ref{prop:square-determinant}. A common finite twist makes its value a specified primitive fourth root. Thus the special tangent exponents satisfy $c_i\equiv1\pmod4$ in $\Z/24$.

A one-slope polygon has a scalar leading-coordinate action on every invariant residual line. If it is reducible, the tame scalar belongs to $\F_5$, and its determinant is square. Thus both polygons have two slopes. Their first-level modules are rational over the common tame good field $L/K$ of degree 24. Lemma~\ref{lem:one-prime-parameter} forces $B_i=25k_i$, with $1\le k_i\le3$. If the high lines agreed, their common characters and determinant would give $k_1\equiv k_2\pmod4$, hence $k_1=k_2$ and equal special characters. The high lines are therefore distinct.

A high-line character takes values in $\F_5^\times$, so $c_i-B_i$ is divisible by six. Since $c_i$ is odd, $k_i$ is odd. Thus the parameters must be 25 and 75. For $B=25$ the possible exponents are $c=1,13$; for $B=75$ they are $c=9,21$. Equality of the exchanged residual character pairs leaves the two rows
\[
 (25,1;75,9),\qquad(25,13;75,21).
\]
A common quadratic twist exchanges them. Put the high lines at $\F_5e_1$ and $\F_5e_2$, in the parameter order 25,75. The two tame residual exponents are $0,6$, up to adding twelve to both.

Write $F_2=L(M_i[25])$, $F_3=L(M_i[125])$ and let $W_1,W_2$ be their congruence layers. We first prove that evaluation $W_1e_i\to\F_5^2$ is onto. A second-level lift of $e_i$ has, in the two models, valuations
\[
          (19/5,3/5)\quad(i=1),\qquad(1/5,9/5)\quad(i=2)
\]
over $L$. Its common field has degree five or 25. At degree five, both normalized values are prime to five. A translation line can be the first high line, the second high line, or a line different from both. The two break values for these three possibilities would be
\[
 \begin{array}{c|ccc}
 i=1&(456,72)&(6,222)&(6,72)\\
 i=2&(474,66)&(24,216)&(24,66).
 \end{array}
\]
For example, in the first row the leading coordinate values are 19 and three; the high/low translating values are 475,25 and 225,75. Subtraction gives the displayed pairs. No pair agrees, contradicting the exact prime-valuation break equality. Thus both point fields $E_i$ have degree 25, and both evaluations are onto.

Tame projectors separate the diagonal plane and the two off-diagonal axes, whose weights are $0,6,18$. The two evaluations consequently put both off-diagonal matrix units in $W_1$ and make both diagonal projections nonzero. In particular $r=\dim W_1\ge3$. Since $W_1\subset W_2$, each lift $Q_i$ of order 125 gives a common field $T_i=L(Q_i)$ with $[T_i:E_i]=25$. A prime-to-five tame lift diagonalizes on the torsion module by the coprime factors of $X^{24}-1$; choose $Q_i$ in its two direct summands. This makes all four point fields stable under the tame lift.

Here are the normalized coordinate and displacement rows. The designation in the last column is the line on which the larger displacement occurs.
\[
\begin{array}{c|c|r|r|r|c}
\text{extension}&\text{model}&A&D_0&D_1&H\\\hline
E_1/L&1&95&125&2375&\text{diagonal}\\
E_1/L&2&15&375&1125&\text{off-diagonal}\\
E_2/L&1&5&125&2375&\text{off-diagonal}\\
E_2/L&2&45&375&1125&\text{diagonal}\\\hline
T_1/E_1&1&95&3125&59375&\text{diagonal}\\
T_1/E_1&2&15&9375&28125&\text{off-diagonal}\\
T_2/E_2&1&5&3125&59375&\text{off-diagonal}\\
T_2/E_2&2&45&9375&28125&\text{diagonal}.
\end{array}
\]
The values come directly from the first-level low/high values $(5,95)$ and $(15,45)$: every inverse step divides its value by 25, while the field degree multiplies by 25. The displacements are the translating first-level values multiplied by the total field degree. The absolute indices are 12000 and 300000, respectively, strictly above every listed displacement.

In each degree 25 extension the translation group has two different tame weights, diagonal and off-diagonal. Its diagonal cyclic quotient break is a positive multiple of 24, and its other quotient break is six or eighteen modulo 24. To see this, on a uniformizer in the cyclic degree-five quotient the tame leading multiplier is $\eta$ with $\eta^5=\zeta$, or $\eta^{125}=\zeta$ at the higher point level. In both cases $\eta=\zeta^5$. Conjugating a wild expansion gives weight $\eta^b=\zeta^{5b}$. Thus the two quotient breaks differ. The group has two lower breaks and its last line is one of the two tame eigenaxes.

The two coordinate rows of each extension have opposite high-displacement lines. At least one of them is not the last ramification line, so Lemma~\ref{lem:opposite-line-conductor} bounds the last upper break by the maximum of $(D_0-A)/5$ for the two rows. We obtain
\[
 c(E_1/L)\le72,\quad c(E_2/L)\le66,\quad
 c(T_1/E_1)\le1872,\quad c(T_2/E_2)\le1866.
\]
Since the two $E_i$ generate $F_2$, its last break is at most 72. The same stabilizer and elementary abelian compositum argument as in Lemma~\ref{lem:point-conductors} gives
\begin{equation}\label{eq:tame-terminal-bound}
                       c(F_3/F_2)\le5^{r-2}\,1872.
\end{equation}
Indeed terminal translations are onto both columns, giving $F_2\cap T_i=E_i$, and $[F_2:E_i]=5^{r-2}$. Upper subgroup conversion bounds the two terminal point quotients by that degree times their original conductors; they generate $F_3/F_2$.

If the first-layer trace vanishes, $\mu_{25}\subset K$ and $L(\mu_{125})/L$ is cyclic of degree five with break 576, by the Cohen-field calculation in Section~\ref{sec:irreducible}. Its break exceeds 72, so its terminal restriction has conductor greater than $5^r(576-72)$, contradicting~\eqref{eq:tame-terminal-bound}, since $25(576-72)=12600>1872$.

Otherwise $\mu_{25}\not\subset K$. Theorem~\ref{thm:cyclotomic-exact} gives $\delta_2\ge20$, so the second break over $L$ is at least $480$. It is greater than 72. The cyclotomic field meets the elementary abelian $F_2$ in its first layer, and its nontrivial terminal restriction has conductor greater than $5^r(480-72)$. This again contradicts~\eqref{eq:tame-terminal-bound}, since
\[
                         25(480-72)=10200>1872.
\]
Both determinant cases are impossible, proving the theorem.
\end{proof}
